% ================= A 卷 =================
% 难度定位：基础巩固（覆盖全部章节最低掌握线）
\juantou{A}
\defenBiaoFour{20}{32}{30}{18}

\dati{名词解释}{每小题 5 分，共 20 分}

\begin{ti}
\xiaoti{1.（5 分）物质}
\kongbai{2.6cm}
\end{ti}

\begin{ti}
\xiaoti{2.（5 分）度}
\kongbai{2.6cm}
\end{ti}

\begin{ti}
\xiaoti{3.（5 分）实践}
\kongbai{2.6cm}
\end{ti}

\begin{ti}
\xiaoti{4.（5 分）社会存在}
\kongbai{2.6cm}
\end{ti}

\dati{简答题}{每小题 8 分，共 32 分}

\begin{ti}
\xiaoti{1.（8 分）哲学基本问题的内容包括哪两个方面？划分唯物主义与唯心主义、可知论与不可知论的标准各是什么？}
\kongbai{4.4cm}
\end{ti}

\begin{ti}
\xiaoti{2.（8 分）为什么说对立统一规律是唯物辩证法的实质和核心？}
\kongbai{4.4cm}
\end{ti}

\begin{ti}
\xiaoti{3.（8 分）实践对认识的决定作用表现在哪些方面？}
\kongbai{4.4cm}
\end{ti}

\begin{ti}
\xiaoti{4.（8 分）社会意识相对独立性的主要表现有哪些？}
\kongbai{4.4cm}
\end{ti}

\dati{论述题}{每小题 15 分，共 30 分}

\begin{ti}
\xiaoti{1.（15 分）试述量变与质变的辩证关系原理，并说明这一原理对个人学习与成长的指导意义。}
\kongbai{9.5cm}
\end{ti}

\begin{ti}
\xiaoti{2.（15 分）论述真理的绝对性与相对性及其辩证关系，并说明把握这一原理的实践意义。}
\kongbai{9.5cm}
\end{ti}

\dati{材料分析题}{共 18 分，阅读所给材料后回答问题}

\begin{ti}
\xiaoti{1.（18 分）某山区县长期交通闭塞、产业单一，群众收入偏低。近年来，当地在调研本地气候与土壤条件的基础上，发展高山蔬菜与林下中药材种植，同时修路、建冷链、培训农户，逐步形成“种植—加工—销售”链条。五年间，全县农村居民人均可支配收入翻了一番，外出务工人员陆续返乡创业。当地干部总结说：“不是没有资源，是过去没找对路子。”\\

请回答：\\
（1）运用社会存在与社会意识辩证关系的原理，分析该县的脱贫实践。（9 分）\\
（2）说明这一案例如何体现人民群众是历史的创造者。（9 分）}
\kongbai{12.0cm}
\end{ti}
